% =====================================================================================
% preamble.tex — cấu hình chung cho cả 4 báo cáo cá nhân
% Biên dịch bằng XeLaTeX (hoặc Tectonic). Xem README.md.
% =====================================================================================
\documentclass[13pt]{extarticle}

% ---------- Trang và font -----------------------------------------------------------
\usepackage[letterpaper, left=2.5cm, right=2cm, top=2.2cm, bottom=2.2cm, footskip=1cm]{geometry}
\usepackage{fontspec}
% Dùng Times New Roman nếu máy có (Windows/Overleaf có sẵn); nếu không, dùng TeX Gyre Termes
% — bản sao Times có đủ dấu tiếng Việt.
\IfFontExistsTF{Times New Roman}{%
  \setmainfont{Times New Roman}%
}{%
  \setmainfont{texgyretermes}[Extension=.otf, UprightFont=*-regular, BoldFont=*-bold,
    ItalicFont=*-italic, BoldItalicFont=*-bolditalic]%
}
\setmonofont{lmmono10}[Extension=.otf, UprightFont=*-regular, Scale=0.9]
\linespread{1.2}
\setlength{\parindent}{0pt}
\setlength{\parskip}{5pt}
\hyphenpenalty=10000
\exhyphenpenalty=10000
\emergencystretch=3em

% ---------- Gói --------------------------------------------------------------------
\usepackage{xcolor}
\usepackage{graphicx}
\usepackage{array, tabularx, longtable, booktabs, multirow, makecell}
\usepackage{colortbl}
\usepackage{enumitem}
\usepackage{titlesec}
\usepackage{tocloft}
\usepackage{caption}
\usepackage{fancyhdr}
\usepackage{bytefield}
\usepackage{url}
\usepackage{etoolbox}
\usepackage{float}
\usepackage{amssymb}
\usepackage{listings}
\lstset{basicstyle=\ttfamily\scriptsize, backgroundcolor=\color{cGray}, frame=single, rulecolor=\color{cLine},
  columns=fullflexible, keepspaces=true, breaklines=true, xleftmargin=2pt, xrightmargin=2pt,
  extendedchars=true, inputencoding=utf8}
\usepackage{tikz}
\usetikzlibrary{positioning, arrows.meta, calc, fit, backgrounds, shapes.multipart,
  shapes.geometric, matrix, decorations.pathreplacing}
\usepackage[hidelinks]{hyperref}

% ---------- Tên tiếng Việt ----------------------------------------------------------
\renewcommand{\figurename}{Hình}
\renewcommand{\tablename}{Bảng}
\renewcommand{\contentsname}{Mục lục}

% ---------- Màu --------------------------------------------------------------------
\definecolor{tocblue}{HTML}{2F5496}
\definecolor{cBlue}{HTML}{DDE8F3}
\definecolor{cGreen}{HTML}{E3F1E8}
\definecolor{cOrange}{HTML}{FBEFD9}
\definecolor{cPurple}{HTML}{F1E3F3}
\definecolor{cGray}{HTML}{EEF0F2}
\definecolor{cHead}{HTML}{E7E7E7}
\definecolor{cLine}{HTML}{555555}
\definecolor{uiBg}{HTML}{2B2F36}
\definecolor{uiBar}{HTML}{15181D}
\definecolor{uiField}{HTML}{1C1F24}
\definecolor{uiBtn}{HTML}{4A525D}
\definecolor{uiText}{HTML}{E8E8E8}
\definecolor{uiState}{HTML}{9FD49F}
\definecolor{uiRed}{HTML}{D9534F}

% ---------- Tiêu đề mục ------------------------------------------------------------
\titleformat{\section}{\bfseries\fontsize{14}{18}\selectfont}{\thesection.}{0.5em}{}
\titleformat{\subsection}{\bfseries\fontsize{13}{16}\selectfont}{\thesubsection.}{0.5em}{}
\titleformat{\subsubsection}{\bfseries\itshape\fontsize{13}{16}\selectfont}{\thesubsubsection.}{0.5em}{}
\titlespacing*{\section}{0pt}{14pt}{6pt}
\titlespacing*{\subsection}{12pt}{10pt}{4pt}
\titlespacing*{\subsubsection}{24pt}{8pt}{3pt}

% "Phần" lớn: I. NỘI DUNG NHÓM / II. NỘI DUNG CÁ NHÂN — sang trang mới, vào mục lục
\newcommand{\PhanLon}[1]{%
  \clearpage\phantomsection
  \addcontentsline{toc}{part}{#1}%
  {\bfseries\fontsize{15}{19}\selectfont #1\par}\vspace{6pt}}

% ---------- Mục lục ----------------------------------------------------------------
\renewcommand{\cfttoctitlefont}{\color{tocblue}\bfseries\fontsize{15}{19}\selectfont}
\renewcommand{\cftaftertoctitle}{}
\setlength{\cftbeforetoctitleskip}{0pt}
\setlength{\cftaftertoctitleskip}{8pt}
\renewcommand{\cftpartfont}{\bfseries}
\renewcommand{\cftpartpagefont}{\bfseries}
\renewcommand{\cftpartleader}{\cftdotfill{\cftdotsep}}
\setlength{\cftbeforepartskip}{4pt}
\renewcommand{\cftsecfont}{\normalfont}
\renewcommand{\cftsecpagefont}{\normalfont}
\renewcommand{\cftsecleader}{\cftdotfill{\cftdotsep}}
\renewcommand{\cftsecaftersnum}{.}
\setlength{\cftbeforesecskip}{2pt}
\setlength{\cftsecindent}{1.2em}
\setlength{\cftsubsecindent}{3em}
\renewcommand{\cftsubsecaftersnum}{.}
\setlength{\cftsecnumwidth}{1.8em}
\setlength{\cftsubsecnumwidth}{2.6em}
\setcounter{tocdepth}{2}
\setcounter{secnumdepth}{3}

% ---------- Hình và bảng -----------------------------------------------------------
\counterwithin{figure}{section}
\captionsetup{font={small,it}, labelfont={small,it}, labelsep=period, justification=centering}
\setlength{\abovecaptionskip}{4pt}
\renewcommand{\arraystretch}{1.25}
\newcolumntype{L}[1]{>{\raggedright\arraybackslash}p{#1}}
\newcolumntype{Y}{>{\raggedright\arraybackslash}X}
\newcommand{\hd}[1]{\cellcolor{cHead}\textbf{#1}}
\setlist{nosep, leftmargin=1.5em, topsep=2pt}

% Mã nguồn trong văn bản: ngắt dòng được ở / . _ và chữ hoa
\urlstyle{tt}
\DeclareUrlCommand\code{\urlstyle{tt}}
\expandafter\def\expandafter\UrlBreaks\expandafter{\UrlBreaks
  \do\A\do\B\do\C\do\D\do\E\do\F\do\G\do\H\do\I\do\J\do\K\do\L\do\M\do\N\do\O\do\P
  \do\Q\do\R\do\S\do\T\do\U\do\V\do\W\do\X\do\Y\do\Z\do\_\do\.}

% ---------- Số trang ---------------------------------------------------------------
\pagestyle{fancy}
\fancyhf{}
\cfoot{\thepage}
\renewcommand{\headrulewidth}{0pt}
\fancypagestyle{plain}{\fancyhf{}\cfoot{\thepage}\renewcommand{\headrulewidth}{0pt}}

% ---------- Style TikZ dùng chung ---------------------------------------------------
\tikzset{
  every picture/.style={font=\small},
  box/.style={draw=cLine, rounded corners=2pt, align=center, inner sep=4pt, fill=white,
    font=\small, minimum height=0.9cm},
  bxB/.style={box, fill=cBlue}, bxG/.style={box, fill=cGreen}, bxO/.style={box, fill=cOrange},
  bxP/.style={box, fill=cPurple}, bxY/.style={box, fill=cGray},
  dec/.style={draw=cLine, diamond, aspect=2.2, align=center, inner sep=1pt, fill=cOrange, font=\small},
  db/.style={draw=cLine, cylinder, shape border rotate=90, aspect=0.25, fill=cGray, align=center,
    minimum height=1cm, font=\small},
  arr/.style={-{Stealth[length=2.2mm]}, draw=cLine, thick},
  darr/.style={arr, dashed},
  biarr/.style={{Stealth[length=2.2mm]}-{Stealth[length=2.2mm]}, draw=cLine, thick},
  lbl/.style={font=\footnotesize, fill=white, inner sep=1.5pt, align=center},
  grp/.style={draw=cLine, dashed, rounded corners=4pt, inner sep=8pt},
  grpt/.style={font=\footnotesize\bfseries, anchor=south west, inner sep=2pt},
  % trạng thái
  st/.style={draw=cLine, rounded corners=5pt, fill=cBlue, minimum height=0.8cm, minimum width=2.4cm,
    align=center, font=\small},
  init/.style={circle, fill=black, minimum size=7pt, inner sep=0pt},
  final/.style={circle, draw=black, double, double distance=1.5pt, fill=black, minimum size=5pt, inner sep=0pt},
  % lớp UML
  umlclass/.style={draw=cLine, rectangle split, rectangle split parts=2, fill=white,
    rectangle split part fill={cBlue, white}, align=left, font=\footnotesize, inner sep=3pt},
}

% ---------- Sơ đồ tuần tự: macro tự viết để kiểm soát vị trí chữ ----------------------
% Dùng trong tikzpicture:
%   \seqactor{id}{x}{Tên}   ... \seqbegin
%   \seqmsg{từ}{đến}{nhãn}   \seqret{từ}{đến}{nhãn}   \seqself{id}{nhãn}
%   \seqnote{trái}{phải}{số dòng}{nội dung}
%   \seqloopbegin   ... \seqloopend{trái}{phải}{nhãn}
%   \seqend
\newdimen\seqy
\newdimen\seqloopy
\newcommand{\seqstep}{0.72cm}
\newcommand{\seqactor}[3]{%
  \node[draw=cLine, fill=cBlue, rounded corners=2pt, minimum height=0.8cm, align=center, font=\small\bfseries]
    (#1) at (#2,0) {#3};}
\newcommand{\seqbegin}{\global\seqy=-0.55cm}
\newcommand{\seqmsg}[3]{%
  \global\advance\seqy by -\seqstep
  \draw[arr] (#1.center |- 0,\the\seqy) -- node[lbl, above=1pt]{#3} (#2.center |- 0,\the\seqy);}
\newcommand{\seqret}[3]{%
  \global\advance\seqy by -\seqstep
  \draw[darr] (#1.center |- 0,\the\seqy) -- node[lbl, above=1pt]{#3} (#2.center |- 0,\the\seqy);}
\newcommand{\seqlost}[3]{% gói bị mất: mũi tên đứt đoạn dừng giữa chừng với dấu X
  \global\advance\seqy by -\seqstep
  \draw[darr, -] (#1.center |- 0,\the\seqy) -- node[lbl, above=1pt]{#3} ($(#1.center |- 0,\the\seqy)!0.7!(#2.center |- 0,\the\seqy)$)
    node[text=red!70!black, font=\bfseries, inner sep=0pt] {$\times$};}
\newcommand{\seqself}[2]{%
  \global\advance\seqy by -\seqstep
  \draw[arr] (#1.center |- 0,\the\seqy) -- ++(0.6,0) -- ++(0,-0.4) -- ++(-0.6,0);
  \node[lbl, anchor=west, align=left] at ($(#1.center |- 0,\the\seqy)+(0.7,-0.2)$) {#2};
  \global\advance\seqy by -0.3cm}
\newcommand{\seqnote}[4]{%
  \global\advance\seqy by -0.45cm
  \path let \p1=(#1.center), \p2=(#2.center), \n1={max(\x2-\x1+1.2cm, 3.4cm)} in
    node[draw=cLine, fill=cOrange, anchor=north, text width=\n1-8pt, align=center, font=\footnotesize,
      inner sep=3pt] at ($(#1.center |- 0,\the\seqy)!0.5!(#2.center |- 0,\the\seqy)$) {#4};
  \global\advance\seqy by -#3\baselineskip
  \global\advance\seqy by -0.05cm}
\newcommand{\seqloopbegin}{\global\advance\seqy by -0.25cm \global\seqloopy=\seqy}
\newcommand{\seqloopend}[3]{%
  \global\advance\seqy by -0.3cm
  \draw[cLine, thin] ($(#1.center |- 0,\the\seqloopy)+(-0.9,0)$) rectangle ($(#2.center |- 0,\the\seqy)+(0.9,0)$);
  \node[anchor=north west, fill=cGray, draw=cLine, font=\scriptsize\bfseries, inner sep=2pt]
    at ($(#1.center |- 0,\the\seqloopy)+(-0.9,0)$) {#3};}
\newcommand{\seqend}[1]{% #1 = danh sách id, ví dụ {C,M,G}
  \global\advance\seqy by -0.4cm
  \begin{scope}[on background layer]
    \foreach \a in {#1} { \draw[cLine!70, dashed] (\a.south) -- (\a.center |- 0,\the\seqy); }
  \end{scope}}

% ---------- Hộp bước đánh số (dùng cho luồng xử lý) ---------------------------------
% \buoc{tên node}{vị trí}{số}{tiêu đề}{mô tả}{màu}
\newcommand{\buoc}[6]{%
  \node[draw=cLine, rounded corners=2pt, fill=#6, text width=4.3cm, align=left, inner sep=4pt,
    minimum height=1.75cm, anchor=north west] (#1) at #2 {%
    \textbf{#3.~#4}\par\vspace{1pt}{\footnotesize #5}};}

% ---------- Giao diện IMGUI (mô phỏng) ---------------------------------------------
\newcommand{\uinum}[1]{\tikz[baseline=-0.6ex]\node[circle, fill=uiRed, text=white, inner sep=0pt,
  minimum size=10pt, font=\tiny\bfseries]{#1};}

% ---------- Bảng ánh xạ ------------------------------------------------------------
\newenvironment{bangAnhXa}{%
  \small
  \begin{longtable}{|L{4.2cm}|L{6.6cm}|L{4.4cm}|}
  \hline
  \hd{Hình / phần tử trên hình} & \hd{Chức năng (mô tả)} & \hd{Hiện thực trong mã nguồn}\\ \hline
  \endfirsthead
  \hline
  \hd{Hình / phần tử trên hình} & \hd{Chức năng (mô tả)} & \hd{Hiện thực trong mã nguồn}\\ \hline
  \endhead
  \hline
  \endfoot
}{\end{longtable}}
\newcommand{\hang}[3]{#1 & #2 & #3 \\ \hline}
